% L5a student companion deck.
%
% The frame order mirrors the lecture notebook so students can take notes while
% the instructor works through the notebook. The disclaimer is intentionally
% slide 2, by instructor preference.
\documentclass[aspectratio=169,11pt]{beamer}
\usepackage{vnslides}

\newcommand{\Pmeasure}{\mathbb{P}}
\newcommand{\E}{\mathbb{E}}
\newcommand{\Var}{\operatorname{Var}}
\newcommand{\Cov}{\operatorname{Cov}}
\newcommand{\gbar}{\bar g}
\newcommand{\bA}{\mathbf{A}}
\newcommand{\bC}{\mathbf{C}}
\newcommand{\bG}{\mathbf{G}}
\newcommand{\bZ}{\mathbf{Z}}
\newcommand{\bw}{\mathbf{w}}
\newcommand{\bg}{\mathbf{g}}
\newcommand{\Sg}{\hat{\boldsymbol{\Sigma}}_{g}}
\newcommand{\Ch}{\hat{\mathbf{C}}}

\title{\texorpdfstring{{\vnheavy L5a}}{L5a}: Single Asset Geometric Brownian Motion and NPV}
\subtitle{CHEME 5660 Cornell University}
\author{Copyright Jeffrey Varner 2026. All Rights Reserved}

\newcommand{\mug}{\mu_g}
\begin{document}
\vntitlepage
\begin{frame}{Disclaimer and Risks}
  \vspace{0.6em}
  \textbf{\ink{Training only.}} This content is for training and informational
  purposes only. The teaching team does not make, give, or endorse any offer or
  solicitation to buy or sell securities or derivative products, or provide
  investment or trading advice or strategy.

  \vspace{1.1em}
  \textbf{\ink{Risk.}} Trading involves risk and is generally not recommended for
  those with limited resources, experience, or a low risk tolerance. Past
  performance does not guarantee future success.

  \vspace{1.1em}
  \textbf{\ink{Responsibility.}} You are responsible for your investment decisions
  based on your financial circumstances, objectives, risk tolerance, and
  liquidity needs.
\end{frame}

\begin{frame}{Objectives for Today}
By the end of this lecture, you will be able to:
\vspace{0.8em}
\begin{itemize}\setlength{\itemsep}{0.8em}
\item \hilite{Interpret out-of-sample predictions}: explain how historical
mean growth and volatility determine the distribution of later prices.
\item \hilite{Calculate an NPV target probability}: connect the discounted
trade payoff to the probability of exceeding a target at the sale date.
\item \hilite{Compare predictions with observations}: use the examples to
assess price predictions and examine whether updating improves forecasts.
\end{itemize}
\end{frame}

\begin{frame}{Lectures and Examples}
\textbf{\ink{Example:}} \href{https://github.com/varnerlab/CHEME-5660-CourseRepository-Fall-2026/blob/main/lectures/week-5/L5a/CHEME-5660-L5a-Example-OOS-SAGBM-Fall-2026.ipynb}{Out-of-Sample Single-Asset GBM}

How well do historical estimates describe later prices? Simulate 2025 prices,
construct prediction bands, and compare observed coverage across assets.

\vspace{0.8em}
\textbf{\ink{Example:}} \href{https://github.com/varnerlab/CHEME-5660-CourseRepository-Fall-2026/blob/main/lectures/week-5/L5a/CHEME-5660-L5a-Example-GBM-NPV-TradeRule-Fall-2026.ipynb}{GBM Net Present Value Trade Rule}

How likely is a trade to exceed a target scaled NPV? Calculate the probability,
check it at the median, and examine how it changes with the target.

\vspace{0.8em}
\textbf{\ink{Example:}} \href{https://github.com/varnerlab/CHEME-5660-CourseRepository-Fall-2026/blob/main/lectures/week-5/L5a/CHEME-5660-L5a-Example-EMA-SAGBM-Fall-2026.ipynb}{Updating GBM Parameters with EMA}

Can recent observations improve the forecasts? Update mean growth and volatility,
then compare the resulting forecasts with the frozen model.
\end{frame}

\begin{frame}{Company Profile: Renaissance Technologies}
Renaissance Technologies (RenTec) is a quantitative investment management
firm founded in 1982. Its investment work uses mathematical and statistical
methods, with a data-driven research approach across its funds.

\vspace{0.7em}
Its researchers include people trained in mathematics, physics, computer
science, and related fields.

\vspace{0.7em}
\hilite{Connection to today:} State the assumptions, record a forecast,
and evaluate it against later observations. Our GBM example is a teaching
model, not a description of RenTec's proprietary methods.
\slidenote{Sources:}{\href{https://www.rentec.com/Home.action?about=true}{RenTec overview};
\href{https://www.rentec.com/Home.action?index=true}{firm description}.}
\note{[Sources] https://www.rentec.com/Home.action?about=true ; https://www.rentec.com/Home.action?index=true [Sources]}
\end{frame}

\begin{frame}{Concept Review: Single Asset GBM Model}
  Recall that geometric Brownian motion (GBM) describes share price changes
  through drift and random fluctuations, both proportional to the current price:
  \[
    \frac{dS(t)}{S(t)}
    =\underbrace{\mu\,dt}_{\text{drift}}
    +\underbrace{\sigma\,dW(t)}_{\text{random fluctuations}}.
  \]

  Here, $S(t)$ is the share price (dollars per share), with initial price
  $S(0)=S_0>0$, and $dW(t)$ is a Wiener-process increment.

  \vspace{0.6em}
  With time measured in years, the constant parameters are:
  \begin{itemize}\setlength{\itemsep}{0.65em}
    \item \hilite{Drift $\mu\in\mathbb R$}: the expected proportional price
      change per unit time (year$^{-1}$).
    \item \hilite{Volatility $\sigma>0$}: sets the size of the random
      fluctuations (year$^{-1/2}$).
  \end{itemize}
\end{frame}

\begin{frame}{Concept Review: Growth and Volatility}
  Over an observation interval $\Delta t>0$ years, the growth rate is:
  \[
    g_j=\frac{1}{\Delta t}\ln\!\left(\frac{S_{t_j}}{S_{t_{j-1}}}\right),
    \qquad t_j=j\Delta t.
  \]

  Under GBM, its mean and standard deviation are:
  \[
    \underbrace{\mu_g=\E[g_j]=\mu-\frac{\sigma^2}{2}}_{\text{mean growth rate}},
    \qquad
    \underbrace{\sqrt{\Var(g_j)}=\frac{\sigma}{\sqrt{\Delta t}}}_{\text{growth-rate standard deviation}}.
  \]
  Drift $\mu$ and mean growth rate $\mu_g$ both have units year$^{-1}$,
  but differ by $\sigma^2/2$.

  \vspace{0.5em}
  From the sample growth-rate standard deviation $\sigma_g$, we estimate
  volatility as:
  \[
    \hat\sigma=\sigma_g\sqrt{\Delta t}.
  \]
\end{frame}

\begin{frame}{Concept Review: The One-Step Transition}
  For a grid of $N$ steps, $t_j=j\Delta t$, the exact transition from L4b is:
  \[
    \boxed{S_{t_j}=S_{t_{j-1}}\exp\!\Bigl[\mu_g\Delta t+
      \sigma\sqrt{\Delta t}\,Z_j\Bigr],\qquad j=1,\ldots,N.}
  \]
  Here, $Z_1,\ldots,Z_N$ are independent standard normal draws, one per step.

  \vspace{0.7em}
  For a single horizon $T>0$, we can sample the terminal price directly:
  \[
    S_T=S_0\exp\!\Bigl[\mu_gT+\sigma\sqrt{T}\,Z\Bigr],
    \qquad Z\sim\mathcal N(0,1).
  \]

  \vspace{0.5em}
  Repeated one-step draws generate a price path. The fixed-horizon form
  samples the price at the selected time $T$.
\end{frame}

\begin{frame}{Example: Out-of-Sample Price Predictions}
\textbf{\ink{Example:}}
  \href{https://github.com/varnerlab/CHEME-5660-CourseRepository-Fall-2026/blob/main/lectures/week-5/L5a/CHEME-5660-L5a-Example-OOS-SAGBM-Fall-2026.ipynb}{Test predictions against later observations}

  \vspace{0.7em}
  How well does a model fitted to historical prices describe prices
  observed afterward?

  \vspace{0.5em}
  \begin{itemize}\setlength{\itemsep}{0.6em}
    \item Use the mean growth rate and volatility estimated from 2014--2024 data.
    \item Simulate prices for 2025 and compare prediction bands with
      observed prices.
    \item Measure the fraction of dates with observed prices inside each band,
      first for one firm and then across the dataset.
  \end{itemize}

  \slidenote{Next question:}{how can this price distribution tell us the probability of exceeding a target return at a scheduled sale?}
\end{frame}

\begin{frame}{GBM Trade Rule}
  Consider a scheduled trade: buy $n_0>0$ shares at $S_0$
  at time $0$ and sell at $S_T$ at time $T=N\Delta t$. Assume no dividends
  or trading costs.

  \vspace{0.4em}
  Let $g_y$ be the constant, continuously compounded benchmark growth rate
  (year$^{-1}$) corresponding to yield $y$. The NPV is given by:
  \[
    \operatorname{NPV}(g_y,T)
    =\underbrace{-n_0S_0}_{\text{purchase today}}
    +\underbrace{n_0S_Te^{-g_yT}}_{\text{discounted sale}}.
  \]

  Dividing by the initial investment gives the \hilite{scaled NPV}, our
  discounted fractional return:
  \[
    \rho_T=\frac{\operatorname{NPV}(g_y,T)}{n_0S_0}
    =\Bigl(\frac{S_T}{S_0}\Bigr)e^{-g_yT}-1.
  \]

  We now use GBM to describe the uncertain sale price $S_T$.
\end{frame}

\begin{frame}{Substitute the GBM Solution}
  Using mean growth rate $\mug=\mu-\sigma^2/2$, substitute the GBM price into
  the scaled NPV:
  \[
    \boxed{\rho_T=\exp\!\left[(\mug-g_y)T+\sigma\sqrt T\,Z\right]-1,
      \qquad Z\sim\mathcal N(0,1).}
  \]
  Discounting subtracts the benchmark growth rate $g_y$ from the mean growth
  rate $\mug$, while leaving the random shock unchanged. The quantity
  $1+\rho_T$ is lognormal, so:
  \[
    \ln(1+\rho_T)\sim\mathcal N\!\left((\mug-g_y)T,\sigma^2T\right).
  \]
  The scaled NPV is a lognormal variable shifted down by one, with
  $\rho_T>-1$.

  \vspace{0.5em}
  The benchmark affects the return through $g_yT$. When $|g_y|T\ll1$,
  discounting is small and $\rho_T\approx S_T/S_0-1$. We use the exact
  discounted expression to calculate the target probability.
\end{frame}

\begin{frame}{Terminal Target Probability}
  Choose a target scaled NPV $\rho_\star>-1$. Adding one and taking the
  logarithm preserves the strict target condition:
  \[
    \rho_T>\rho_\star
    \quad\Longleftrightarrow\quad
    \ln(1+\rho_T)>\ln(1+\rho_\star).
  \]
  Substituting the GBM expression gives:
  \[
    (\mug-g_y)T+\sigma\sqrt T\,Z>\ln(1+\rho_\star).
  \]
  Subtracting $(\mug-g_y)T$ and dividing by $\sigma\sqrt T>0$ gives a
  threshold for the standard normal random variable $Z$:
  \[
    Z>z_\star,\qquad
    \boxed{z_\star=\frac{\ln(1+\rho_\star)-(\mug-g_y)T}{\sigma\sqrt T}.}
  \]
  Thus $Z>z_\star$ means the trade exceeds its target at the scheduled sale.
  Selling when a boundary is first reached requires a separate calculation.
\end{frame}

\begin{frame}{The Closed-Form Target Probability}
  The standard normal cumulative distribution function $\Phi(z_\star)$ gives
  the probability that $Z\le z_\star$. Subtracting from one gives the
  probability that $Z>z_\star$:
  \[
    \boxed{\Pmeasure(\rho_T>\rho_\star)
      =1-\Phi(z_\star)
      =1-\Phi\!\left(
        \frac{\ln(1+\rho_\star)-(\mug-g_y)T}{\sigma\sqrt T}\right).}
  \]

  \vspace{0.4em}
  \begin{itemize}\setlength{\itemsep}{0.5em}
    \item For a chosen target, we evaluate $z_\star$ from estimated mean growth
      and volatility, the holding period, and benchmark growth rate.
    \item The formula applies for $T>0$, $\sigma>0$, and $\rho_\star>-1$.
    \item A target of $-1$ or below is exceeded with probability one because
      the scaled NPV always exceeds $-1$ under GBM.
  \end{itemize}
  \slidenote{Example:}{\href{https://github.com/varnerlab/CHEME-5660-CourseRepository-Fall-2026/blob/main/lectures/week-5/L5a/CHEME-5660-L5a-Example-GBM-NPV-TradeRule-Fall-2026.ipynb}{Calculate and verify the NPV target probability}}
\end{frame}

\begin{frame}{Change the Benchmark}
The target scaled NPV corresponds to a required sale price:
\[
K=S_0(1+\rho_\star)e^{g_yT}.
\]
With $\rho_\star=0$, the trade succeeds if it beats the benchmark.

\vspace{0.6em}
\begin{itemize}\setlength{\itemsep}{0.55em}
\item $g_y=0$: did the price increase?
\item Positive $g_y$: did it increase enough over the holding period?
\item Increasing $g_y$ raises the sale threshold and lowers its probability.
\end{itemize}
\vspace{0.5em}
The same prices can be assessed against different assumed growth benchmarks.
Choose the benchmark before examining the outcome.
\end{frame}

\begin{frame}{Example: GBM Net Present Value Trade Rule}
\href{https://github.com/varnerlab/CHEME-5660-CourseRepository-Fall-2026/blob/main/lectures/week-5/L5a/CHEME-5660-L5a-Example-GBM-NPV-TradeRule-Fall-2026.ipynb}{Evaluate an NPV-based GBM trade rule}

\vspace{0.8em}
How likely is a stock trade to exceed our target scaled NPV at a scheduled sale?

\vspace{0.6em}
We calculate this probability using the selected firm's mean growth and
volatility, the holding period, and the benchmark growth rate.

\vspace{0.6em}
We check the calculation at the median scaled NPV, which must be exceeded
with probability one half, then vary the target to examine the probability curve.
\end{frame}

\begin{frame}{Example: Updating the Parameter Estimates}
\href{https://github.com/varnerlab/CHEME-5660-CourseRepository-Fall-2026/blob/main/lectures/week-5/L5a/CHEME-5660-L5a-Example-EMA-SAGBM-Fall-2026.ipynb}{Compare frozen and updated estimates}

\vspace{0.8em}
Can giving recent observations more weight improve our forecasts?

\vspace{0.7em}
We retain the original purchase and update its target probability each day
for a sale at the end of the chosen forward window. The example compares frozen
parameters with exponentially weighted estimates of volatility alone or
both mean growth and volatility.

\vspace{0.7em}
We compare the forecasts on the same observed outcomes to examine whether
the updates improve predictions or add noise.
\end{frame}

\begin{frame}{Optional Advanced Material}
The L4b examples develop two related questions:
\vspace{0.7em}
\begin{itemize}\setlength{\itemsep}{0.9em}
\item \hilite{Uncertainty in mean growth}: why can an estimate remain
uncertain even with thousands of observations? Examine the observation
period, sampling frequency, and their effects on target probabilities.
\item \hilite{Monte Carlo versus the closed form}: how many paths are
needed to estimate a probability accurately? Compare simulation error
with the analytical result and introduce variance reduction techniques.
\end{itemize}
\slidenote{Examples:}{\href{https://github.com/varnerlab/CHEME-5660-CourseRepository-Fall-2026/blob/main/lectures/week-4/L4b/advanced/drift-uncertainty/CHEME-5660-L4b-Advanced-DriftUncertainty-Fall-2026.ipynb}{Mean-growth uncertainty};
\href{https://github.com/varnerlab/CHEME-5660-CourseRepository-Fall-2026/blob/main/lectures/week-4/L4b/advanced/monte-carlo/CHEME-5660-L4b-Advanced-MonteCarlo-TargetProbability-Fall-2026.ipynb}{Monte Carlo target probabilities}.}
\end{frame}

\begin{frame}{Summary}
We connected the single-asset GBM model with discounted trade payoffs and
observed outcomes.
\vspace{0.6em}
\begin{itemize}\setlength{\itemsep}{0.75em}
\item \hilite{Out-of-sample predictions}: we used historical estimates to
describe the distribution of later prices.
\item \hilite{NPV target probabilities}: we derived the chance of exceeding
a target and connected it to the chosen benchmark growth rate.
\item \hilite{Predictions and observations}: we checked price predictions
and examined whether updated estimates improve target probabilities.
\end{itemize}
\slidenote{Next:}{L5b extends the price model to correlated assets and portfolio weights.}
\end{frame}
\end{document}
